% Clase 1 — Aplicaciones en IA: Demos
\section{Aplicaciones en IA}

\subsection{Machine Learning}

Paradigma donde los sistemas \textbf{aprenden de datos} en lugar de ser programados explícitamente.

\begin{itemize}
    \item \textbf{Aprendizaje supervisado}: clasificación, regresión.
    \item \textbf{Aprendizaje no supervisado}: clustering, reducción de dimensionalidad.
    \item \textbf{Aprendizaje por refuerzo (RL)}: aprender mediante interacción con el entorno y recompensas.
\end{itemize}

\subsection{Deep Learning}

Redes neuronales profundas con múltiples capas de abstracción.
\begin{itemize}
    \item CNN (Convolutional Neural Networks): visión computacional.
    \item RNN / Transformers: procesamiento de lenguaje natural (NLP).
    \item LLMs: modelos de lenguaje de gran escala (ej. GPT, Claude).
\end{itemize}

\subsection{Áreas de Aplicación}

\begin{itemize}
    \item \textbf{NLP}: traducción automática, chatbots, análisis de sentimientos.
    \item \textbf{Visión computacional}: detección de objetos, reconocimiento facial, diagnóstico médico por imágenes.
    \item \textbf{Robótica}: percepción, planeación de movimiento, manipulación.
\end{itemize}

% Agrega aquí resúmenes de las demos vistas en clase
